\section{Harmonic oscillators}\label{sec:harmonic}

\motiv{\vetnew{The bounds so far are for bounded local generators: spins and open systems. Harmonic
lattice bosons have an unbounded Hamiltonian and an infinite-dimensional algebra, so the generator
of their dynamics has no finite strength and the path-sum bound of
Chapter~\ref{sec:master} does not apply to them. What replaces it is a
different mechanism, special to quadratic Hamiltonians: their dynamics moves the Weyl operators
along a linear flow of the classical phase space, $\alpha_t(W(f))=W(T_tf)$, and the Weyl relation
turns the commutator of two Weyl operators into a phase, so that the whole question reduces to
the entries of the classical propagator $T_t=e^{tK}$.

We bound these entries by the
exponential of the entrywise absolute value $\abs K$, a nonnegative matrix in the role
of the site matrix of Chapter~\ref{sec:closed}: the range form $(vt)^r/r!\,e^{vt}$ and the light
cone come over unchanged, with the classical generator in place of the interaction.}}

\subsection{The lattice phase space and the Weyl operators}

\vetnew{Throughout, $\Lambda$ is a finite set of sites. The phase space is
$V=\Lambda\to\R^2$, a position coordinate $f^q_x$ and a momentum coordinate $f^p_x$ at every
site, equivalently a function on the coordinates $\Lambda'=\Lambda\times\{q,p\}$. For a
coordinate $a=(x,q)$ write $\bar a=(x,p)$ and conversely, and $\site(a)=x$. The symplectic form
is
\begin{equation}\label{eq:sigma}
\sigma(f,g)=\sum_{x\in\Lambda}\bigl(f^q_xg^p_x-f^p_xg^q_x\bigr)=f\cdot\Omega g,
\qquad \Omega_{(x,q),(x,p)}=1,\ \Omega_{(x,p),(x,q)}=-1,
\end{equation}
all other entries of $\Omega$ being zero; it is bilinear and alternating,
$\sigma(f,f)=0$. The support of $f$ is the set of sites where a coordinate of $f$ is nonzero, and
for $X\subseteq\Lambda$ we write $X'=X\times\{q,p\}$ and $\norm f_1=\sum_a\abs{f_a}$.}

\begin{definition}[Weyl system, CCR algebra]\label{def:weylsystem}
\vetnew{A \emph{Weyl system} for $\sigma$ on a unital C*-algebra $\cA$ is a map $W:V\to\cA$ with
every $W(f)$ unitary, the \emph{Weyl relation}
\begin{equation}\label{eq:weylrel}
W(f)W(g)=e^{-\ii\sigma(f,g)/2}\,W(f+g),
\end{equation}
and the $W(f)$ generating $\cA$ as a C*-algebra. The \emph{CCR algebra} of $\sigma$ is the
C*-algebra generated by the concrete Weyl operators on $\ell^2(V)$, $V$ taken with the discrete
topology (so $\ell^2(V)$ is not separable); it is the $\cA$ of this chapter, with its Weyl
system.}
\end{definition}

\vetnew{From \eqref{eq:weylrel}, $W(0)=1$ (as $W(0)=W(0)^2$ with $W(0)$ invertible) and
$W(f)^*=W(-f)$ (as $W(f)W(-f)=W(0)$). In the oscillator picture $W(f)=\exp(\ii\sum_x(f^q_xq_x+f^p_xp_x))$
with $[q_x,p_y]=\ii\delta_{xy}$, and the phase in \eqref{eq:weylrel} is that of the product
formula for the exponentials of two operators with a central commutator; the C*-algebraic definition uses the relation and not the
picture.}

\begin{definition}[the bosonic net]\label{def:bosonicnet}
\vetnew{For $X\subseteq\Lambda$, $\cA(X)$ is the C*-subalgebra of $\cA$ generated by the $W(f)$
with $\supp f\subseteq X$.}
\end{definition}

\begin{proposition}[the form and the net]\label{prop:bosonicnet}
\vetnew{\begin{enumerate}[label=(\alph*),nosep]
\item If $\supp f\cap\supp g=\emptyset$ then $\sigma(f,g)=0$.
\item $X\mapsto\cA(X)$ is a net, causal for the lattice index: for $X\cap Y=\emptyset$, every
element of $\cA(X)$ commutes with every element of $\cA(Y)$.
\end{enumerate}}
\end{proposition}
\begin{proof}
\vetnew{\begin{enumerate}[label=(\alph*),nosep]
\item Every summand of \eqref{eq:sigma} has a factor from a site outside one of the supports.
\item Isotony is clear. For $\supp f\subseteq X$, $\supp g\subseteq Y$,
Theorem~\ref{thm:weylcommutator} below with (a) gives $W(f)W(g)=W(g)W(f)$. The commutant of a set
closed under adjoints is a C*-subalgebra, and $\{W(g):\supp g\subseteq Y\}$ is closed under
adjoints ($W(g)^*=W(-g)$); so this commutant contains the generators of $\cA(X)$, hence
$\cA(X)$, and then the commutant of $\cA(X)$ contains the generators of $\cA(Y)$, hence
$\cA(Y)$.\qedhere
\end{enumerate}}
\end{proof}

\begin{theorem}[the Weyl commutator]\label{thm:weylcommutator}
\vetnew{For all $f,g\in V$,
\begin{equation}\label{eq:weylcomm}
\comm{W(f)}{W(g)}=\bigl(e^{-\ii\sigma(f,g)}-1\bigr)\,W(g)W(f),\qquad
\norm{\comm{W(f)}{W(g)}}=2\,\bigl|\sin\tfrac{\sigma(f,g)}2\bigr|\le\abs{\sigma(f,g)}.
\end{equation}}
\end{theorem}
\begin{proof}
\vetnew{By \eqref{eq:weylrel} twice, $W(f)W(g)=e^{-\ii\sigma(f,g)/2}W(f+g)$ and
$W(g)W(f)=e^{\ii\sigma(f,g)/2}W(f+g)$, so $W(f)W(g)=e^{-\ii\sigma(f,g)}W(g)W(f)$, which is the
identity. $W(g)W(f)$ is unitary, so the norm is $\abs{e^{-\ii s}-1}=2\abs{\sin(s/2)}$ with
$s=\sigma(f,g)$, and $\abs{\sin(s/2)}\le\abs s/2$.}
\end{proof}

\vetnew{The commutator of two Weyl operators is thus a number times a unitary, and its size is
a function of $\sigma(f,g)$ alone: this is what makes the quadratic case elementary.}

\subsection{Quadratic generators and the flow}

\begin{definition}[quadratic generator, flow, dynamics]\label{def:quadgen}
\vetnew{A \emph{quadratic generator} is a real matrix $K$ on $\Lambda'$ with
\begin{equation}\label{eq:quadgen}
K^{\mathsf T}\Omega+\Omega K=0,\qquad\text{equivalently}\quad \sigma(Ku,v)+\sigma(u,Kv)=0.
\end{equation}
Its \emph{flow} is $T_t=e^{tK}$, $t\in\R$. The \emph{dynamics} $\alpha_t$ is the
$*$-automorphism of $\cA$ induced by the symplectic linear bijection $T_t$ of
Proposition~\ref{prop:flow}, the one with
\begin{equation}\label{eq:alphaweyl}
\alpha_t(W(f))=W(T_tf).
\end{equation}}
\end{definition}

\vetnew{The condition \eqref{eq:quadgen} says that $K$ is a Hamiltonian matrix. For the quadratic
Hamiltonian $H=\tfrac12\xi^{\mathsf T}M\xi$ in the coordinates $\xi=(q,p)$, $M$ symmetric,
Hamilton's equations are $\dot\xi=JM\xi$ with $J=\Omega$, and in the Heisenberg picture
$\xi(t)=e^{tJM}\xi$. The Weyl operator $W(f)=e^{\ii f\cdot\xi}$ therefore evolves into
$e^{\ii f\cdot\xi(t)}=e^{\ii(e^{tK}f)\cdot\xi}$ with $K=(JM)^{\mathsf T}=-MJ$: the \emph{label} $f$
moves by the transpose of Hamilton's matrix, which is again Hamiltonian. The generator of this
chapter is that transpose, and \eqref{eq:alphaweyl} is its defining property; the identification
with $e^{\ii tH}\,\cdot\,e^{-\ii tH}$ on Fock space is not part of what is proved.}

\begin{proposition}[the flow]\label{prop:flow}
\vetnew{For a quadratic generator $K$: $T_0=1$, $T_sT_t=T_{s+t}$, $T_t^{-1}=T_{-t}$, and $T_t$ is
symplectic, $T_t^{\mathsf T}\Omega T_t=\Omega$, that is $\sigma(T_tu,T_tv)=\sigma(u,v)$.
Consequently $\alpha_t$ exists and is a $*$-automorphism with \eqref{eq:alphaweyl}.}
\end{proposition}
\begin{proof}
\vetnew{The group law is the exponential of commuting matrices. For symplecticity,
\eqref{eq:quadgen} with $\Omega^{-1}=-\Omega$ (from $\Omega^2=-1$) reads
$K^{\mathsf T}=\Omega(-K)\Omega^{-1}$, so $(e^{tK})^{\mathsf T}=e^{tK^{\mathsf T}}=\Omega e^{-tK}\Omega^{-1}$
and $T_t^{\mathsf T}\Omega T_t=\Omega e^{-tK}e^{tK}=\Omega$. A symplectic linear bijection
$T$ of the phase space induces the automorphism $W(f)\mapsto W(Tf)$ of the CCR algebra: the
relabelling $\psi\mapsto\psi\circ T^{-1}$ of $\ell^2(V)$ is unitary, and conjugation by it sends
$W(f)$ to $W(Tf)$, hence maps the C*-algebra generated by the Weyl operators onto itself; the
automorphisms compose as the maps do. The group $t\mapsto\alpha_t$ is not continuous in norm
(Weyl operators with distinct labels are at distance $2$); \emph{dynamics} here means the
automorphism group.}
\end{proof}

\vetnew{One checks that $\alpha_s\circ\alpha_t=\alpha_{s+t}$: both sides are $*$-automorphisms,
they agree on the Weyl operators, $W(f)\mapsto W(T_sT_tf)=W(T_{s+t}f)$, and a $*$-automorphism is
determined by its values on the Weyl operators, which generate $\cA$.}

\subsection{The entrywise bound}

\vetnew{For a real matrix $K$ over any finite index set write $\abs K$ for the matrix of absolute
values, $\abs K_{ba}=\abs{K_{ba}}$.}

\begin{lemma}[powers]\label{lem:powers}
\vetnew{$\abs{(K^n)_{ba}}\le(\abs K^n)_{ba}$ for all $n$, and $(\abs K^n)_{ba}\ge0$.}
\end{lemma}
\begin{proof}
\vetnew{Induction: $\abs{(K^{n+1})_{ba}}=\abs{\sum_c(K^n)_{bc}K_{ca}}\le\sum_c(\abs K^n)_{bc}\abs K_{ca}$.}
\end{proof}

\begin{theorem}[entrywise exponential bound]\label{thm:entrywise}
\vetnew{For a real matrix $K$ over a finite index set and $t\in\R$,
\begin{equation}\label{eq:entrywise}
\abs{(e^{tK})_{ba}}\le\bigl(e^{\abs t\,\abs K}\bigr)_{ba}=\sum_{n\ge0}\frac{\abs t^n}{n!}(\abs K^n)_{ba}\ \ge0 .
\end{equation}}
\end{theorem}
\begin{proof}
\vetnew{The exponential series converges in norm, hence entrywise, on both sides; termwise
$\abs{t^n(K^n)_{ba}/n!}\le\abs t^n(\abs K^n)_{ba}/n!$ by Lemma~\ref{lem:powers}, and the
triangle inequality for absolutely convergent series gives the bound.}
\end{proof}

\begin{definition}[range, row sum]\label{def:rangerowsum}
\vetnew{Let $d$ be a distance on $\Lambda$ ($d(x,x)=0$, the triangle inequality) and
$\site$ a map from the index set to $\Lambda$. $K$ has \emph{range} $R\ge0$ if $K_{ba}=0$ whenever
$d(\site b,\site a)>R$, and \emph{row sum} $v\ge0$ if $\sum_a\abs{K_{ba}}\le v$ for every $b$.}
\end{definition}

\begin{lemma}[range and row sum]\label{lem:rangerowsum}
\vetnew{With range $R$ and row sum $v$: $(\abs K^n)_{ba}=0$ whenever $nR<d(\site b,\site a)$, and
$\sum_a(\abs K^n)_{ba}\le v^n$.}
\end{lemma}
\begin{proof}
\vetnew{Both by induction on $n$. A nonzero term $\abs K_{bc}(\abs K^n)_{ca}$ needs
$d(\site b,\site c)\le R$ and $d(\site c,\site a)\le nR$, hence $d(\site b,\site a)\le(n+1)R$.
Row sums: $\sum_a(\abs K^{n+1})_{ba}=\sum_c\abs K_{bc}\sum_a(\abs K^n)_{ca}\le v^{n+1}$.}
\end{proof}

\begin{corollary}[range form, light cone]\label{cor:rangeform}
\vetnew{With range $R$, row sum $v$, $t\ge0$, and an integer $r\ge0$ with $nR<d(\site b,\site a)$ for
every $n<r$,
\begin{equation}\label{eq:rangeform}
\bigl(e^{t\abs K}\bigr)_{ba}\le\frac{(vt)^r}{r!}\,e^{vt},
\qquad\text{and}\qquad
\bigl(e^{t\abs K}\bigr)_{ba}\le e^{-\mu r}\bigl(e^{e^\mu vt}-1\bigr)\quad(\mu\ge0,\ r\ge1).
\end{equation}}
\end{corollary}
\begin{proof}
\vetnew{Terms with $n<r$ vanish by Lemma~\ref{lem:rangerowsum}, and the others are at most
$t^nv^n/n!$ by the row-sum bound, so $(e^{t\abs K})_{ba}\le\sum_{n\ge r}(vt)^n/n!$; this tail is
at most $(vt)^re^{vt}/r!$ by Lemma~\ref{lem:tail}. For the cone, $1\le e^{\mu(n-r)}$ for $n\ge r$
and $\mu\ge0$, so the tail is at most $e^{-\mu r}\sum_{n\ge1}(e^\mu vt)^n/n!$ when $r\ge1$.}
\end{proof}

\vetnew{Both bounds in \eqref{eq:rangeform} estimate the tail $\sum_{n\ge r}(vt)^n/n!$, which is itself
an upper bound on the entry. For $R>0$ and $d=d(\site b,\site a)>0$ the
integer $r=\lceil d/R\rceil$ is admissible; for $R=0$ and $d>0$ every $r$ is.}

\subsection{The harmonic bound}

\begin{theorem}[harmonic Lieb--Robinson bound]\label{thm:harmonicbound}
\vetnew{Let $K$ be a quadratic generator, $\supp f\subseteq X$, $\supp g\subseteq Y$ and
$t\in\R$. Then
\begin{equation}\label{eq:harmonicbound}
\norm{\comm{\alpha_t(W(f))}{W(g)}}\le\abs{\sigma(T_tf,g)}\le
\sum_{b\in Y'}\sum_{a\in X'}\bigl(e^{\abs t\,\abs K}\bigr)_{ba}\,\abs{f_a}\,\abs{g_{\bar b}}.
\end{equation}}
\end{theorem}
\begin{proof}
\vetnew{By \eqref{eq:alphaweyl} and Theorem~\ref{thm:weylcommutator}, the left side is at most
$\abs{\sigma(T_tf,g)}$. By \eqref{eq:sigma}, $\sigma(u,g)=\sum_b\epsilon_bu_bg_{\bar b}$ with
$\epsilon_b=\pm1$, so $\abs{\sigma(u,g)}\le\sum_{b\in Y'}\abs{u_b}\abs{g_{\bar b}}$ when
$\supp g\subseteq Y$; and $\abs{(T_tf)_b}\le\sum_{a\in X'}\abs{(e^{tK})_{ba}}\abs{f_a}$, bounded by
\eqref{eq:entrywise}.}
\end{proof}

\vetnew{We state the bound for Weyl operators, as the source does; it extends to finite
linear combinations $\sum_ic_iW(f_i)$ by the triangle inequality, with $\sum_i\abs{c_i}$ in
place of a norm. A bound for arbitrary elements of $\cA(X)$ and $\cA(Y)$ in terms of their norms
is not claimed.}

\begin{corollary}[closed forms]\label{cor:harmonicclosed}
\vetnew{Let $K$ be a quadratic generator of range $R$ and row sum $v$, with $\site(a)$ the site of
the coordinate $a$ and $d$ a distance on $\Lambda$ as in Definition~\ref{def:rangerowsum}, not
assumed symmetric ($d(y,x)$ is read from the site $y$ of the probe); let $\supp f\subseteq X$,
$\supp g\subseteq Y$, $t\ge0$, and let $r\ge0$ be an integer with $nR<d(y,x)$ for all $y\in Y$,
$x\in X$, $n<r$. Then
\begin{equation}\label{eq:harmonicclosed}
\norm{\comm{\alpha_t(W(f))}{W(g)}}\le\norm f_1\norm g_1\,\frac{(vt)^r}{r!}\,e^{vt},
\qquad
\norm{\comm{\alpha_t(W(f))}{W(g)}}\le\norm f_1\norm g_1\,e^{-\mu r}\bigl(e^{e^\mu vt}-1\bigr)
\end{equation}
for $\mu\ge0$, $r\ge1$.}
\end{corollary}
\begin{proof}
\vetnew{Corollary~\ref{cor:rangeform} in \eqref{eq:harmonicbound}, and
$\sum_{b\in Y'}\sum_{a\in X'}\abs{f_a}\abs{g_{\bar b}}\le\norm f_1\norm g_1$.}
\end{proof}

\vetnew{For $R>0$ and $d(Y,X):=\min_{y\in Y,x\in X}d(y,x)>0$, the integer
$r=\lceil d(Y,X)/R\rceil\ge1$ is admissible in Corollary~\ref{cor:harmonicclosed}, and for $\mu>0$
the second bound, at most $\norm f_1\norm g_1e^{-\mu d(Y,X)/R+e^\mu vt}$, is a light cone of
velocity $e^\mu vR/\mu$ (at $\mu=1$: $e\,vR$) and decay rate $\mu/R$ in the distance between the
supports.}

\begin{proposition}[rescaling]\label{prop:rescaling}
\vetnew{Let $K$ be a quadratic generator, $S$ diagonal with positive entries $s_a$ and
$s_{\bar a}=1/s_a$ (a change of scale $q\to s\,q$, $p\to p/s$ at each site, symplectic), and
$K_S=S^{-1}KS$. Then $K_S$ is a
quadratic generator with flow $T^S_t=S^{-1}T_tS$; $\sigma(T^S_tS^{-1}f,S^{-1}g)=\sigma(T_tf,g)$, so
the commutator of $(K_S,S^{-1}f,S^{-1}g)$ is that of $(K,f,g)$; $\abs{K_S}=S^{-1}\abs KS$ and the
bound \eqref{eq:harmonicbound} is unchanged; the row sum of $\abs{K_S}$ is
$\max_bs_b^{-1}\sum_a\abs{K_{ba}}s_a$, which depends on $S$, and the prefactor of
\eqref{eq:harmonicclosed} becomes $\norm{S^{-1}f}_1\norm{S^{-1}g}_1$.}
\end{proposition}
\begin{proof}
\vetnew{$\sigma(Su,Sv)=\sum_a\epsilon_as_au_as_{\bar a}v_{\bar a}=\sigma(u,v)$ as $s_as_{\bar a}=1$, that is
$S^{\mathsf T}\Omega S=\Omega$, so $S$ and $S^{-1}$ are symplectic and
$\sigma(K_Su,v)=\sigma(KSu,Sv)=-\sigma(Su,KSv)=-\sigma(u,K_Sv)$; conjugation commutes with the
exponential series and with the entrywise absolute value, since $s_b^{-1}s_a>0$; in
\eqref{eq:harmonicbound} the weights $s_b^{-1}s_as_a^{-1}s_{\bar b}^{-1}=1$ cancel.}
\end{proof}

\vetnew{The entrywise bound is thus invariant under the choice of units at each site, but the
range form is not: its velocity is the row sum, which the rescaling can lower.}

\subsection{The harmonic chain}

\begin{definition}[the harmonic chain]\label{def:harmonicchain}
\vetnew{On $\Lambda=\{0,\dots,n-1\}$ with $d(x,y)=\abs{x-y}$, $\omega\ge0$ and $\lambda\ge0$,
\begin{equation}\label{eq:chainH}
H=\tfrac12\sum_x\bigl(p_x^2+\omega^2q_x^2\bigr)+\tfrac12\lambda\sum_{x=0}^{n-2}(q_x-q_{x+1})^2
=\tfrac12\sum_xp_x^2+\tfrac12\,q^{\mathsf T}Aq,
\end{equation}
with $A_{xx}=\omega^2+\lambda\deg x$ ($\deg x$ the number of neighbours of $x$),
$A_{x,x\pm1}=-\lambda$, and $A_{xy}=0$ otherwise. Hamilton's equations are $\dot q=p$,
$\dot p=-Aq$, and the generator of the labels is their transpose,
$(Kf)^q=-Af^p$, $(Kf)^p=f^q$.}
\end{definition}

\begin{proposition}[the chain]\label{prop:harmonicchain}
\vetnew{$K$ is a quadratic generator of range $1$ and, for $\lambda\ge0$, of row sum
$v=\max(1,\omega^2+4\lambda)$.}
\end{proposition}
\begin{proof}
\vetnew{$\sigma(Ku,v)=-\sum_{x,y}A_{xy}u^p_yv^p_x-\sum_xu^q_xv^q_x$ is symmetric in $(u,v)$ because
$A$ is, so $\sigma(Ku,v)=\sigma(Kv,u)=-\sigma(u,Kv)$. The range and the row sums are read off
$A$, with $\deg x\le2$.}
\end{proof}

\vetnew{One checks from $A$ that the row sums are exactly $\omega^2+2\lambda\deg x$ on the position
rows and $1$ on the momentum rows. The rescaling $s_{(x,q)}=c$, $s_{(x,p)}=1/c$ of
Proposition~\ref{prop:rescaling} multiplies the position rows by $c^{-2}$ and the momentum rows by
$c^2$, so the row sums become at most $(\omega^2+4\lambda)/c^2$ and $c^2$, balanced at
$c^4=\omega^2+4\lambda$ with $v=(\omega^2+4\lambda)^{1/2}$ when $\omega^2+4\lambda>0$ (for
$\omega=\lambda=0$ the position rows vanish, the row sum becomes $c^2$, and $v$ can be taken
as small as wanted).}

\begin{example}[the bound against the exact commutator]\label{ex:harmonicchain}
\vetnew{Twelve sites, $\omega=\lambda=1$, so $v=5$; $f$ at the site $0$ and $g$ at the site $y$,
both with coordinates of modulus at most $1$. Over this class the maximum of $\abs{\sigma(T_tf,g)}$
is attained at sign vectors (a bilinear form on a box), and the bound \eqref{eq:harmonicbound} is
the sum of the four entries of $e^{t\abs K}$ between the two sites. Figure~\ref{fig:harmonic}
shows both against $y$. At $t=0.5$, $1$, $2$ the exact quantity falls below $0.1$ from $y=2$,
$3$, $4$ on; a least-squares fit of its logarithm over $y=3,\dots,10$ gives the decay per site
$e^{-6.4y}$, $e^{-5.0y}$, $e^{-3.5y}$, and the same fit of the bound $e^{-6.4y}$, $e^{-5.0y}$,
$e^{-3.7y}$ (locally the decay is faster than exponential, as the figure shows).

The bound
exceeds the exact value by factors between $1.04$ and $1.4$ at $t=0.5$, between $1.1$ and $3.4$
at $t=1$, and between $1.7$ and $29$ at $t=2$, the largest factors at small $y$, where the bound
also exceeds the trivial bound $2$.

The range form \eqref{eq:harmonicclosed}
with $r=y$ and the prefactor $\norm f_1\norm g_1=4$ is trivial for every $y$ shown at $t=1$ and
$t=2$, and up to $y=7$ at $t=0.5$: at $t=1$ it reads
$3\cdot10^3$, $7\cdot10^3$, $1\cdot10^4$ for $y=1,2,3$ against $1.1$, $0.18$, $0.0088$, because
$v=5$ counts every entry of $\abs K$ with its full weight, while the entries of $e^{t\abs K}$
between distant sites are small. With the balanced rescaling, $v=\sqrt5$, the cone of
Corollary~\ref{cor:harmonicclosed} (the paragraph after it, at $R=\mu=1$) has velocity $e\,v=6.1$; the chain's group velocity, the
largest slope of its dispersion $\gamma(k)=(\omega^2+4\lambda\sin^2(k/2))^{1/2}$, is
$(\sqrt5-1)/2=0.62$: the cone of the bound is ten times wider than the propagation it bounds.}
\end{example}

\begin{figure}[htb]
\centering
% Figure: the harmonic bound on the chain of chapter 7, generated by checks/m5_figure.py.
% Transcription (for cheap agents): axes, the site y of the second observable to the right
% (the first sits at site 0) and the commutator quantity upward on a logarithmic scale; per
% time t a solid curve with dots (the exact |sigma(T_t f, g)|, maximized over unit f, g) and
% a dashed curve (the entrywise bound), both falling exponentially in y beyond a front that
% moves with t; the bound lies above the exact curve with the same kind of decay. Curves stop
% where they fall below the bottom of the scale; a legend at the top right names the times
% (solid = exact, dashed = bound).
\begin{tikzpicture}[x=0.5cm,y=1cm,>={Stealth[length=5pt]}]
  \draw[->] (0,0) -- (12.6,0) node[right,font=\small] {$y$};
  \draw[->] (0,0) -- (0,6.18) node[above,font=\small] {$|\sigma(T_tf,g)|$};
  \draw (0.15,0.42) -- (-0.15,0.42) node[left,font=\scriptsize] {$10^{-12}$};
  \draw (0.15,1.68) -- (-0.15,1.68) node[left,font=\scriptsize] {$10^{-9}$};
  \draw (0.15,2.94) -- (-0.15,2.94) node[left,font=\scriptsize] {$10^{-6}$};
  \draw (0.15,4.20) -- (-0.15,4.20) node[left,font=\scriptsize] {$10^{-3}$};
  \draw (0.15,5.46) -- (-0.15,5.46) node[left,font=\scriptsize] {$10^{0}$};
  \draw (1,0.06) -- (1,-0.06) node[below,font=\scriptsize] {1};
  \draw (2,0.06) -- (2,-0.06) node[below,font=\scriptsize] {2};
  \draw (3,0.06) -- (3,-0.06) node[below,font=\scriptsize] {3};
  \draw (4,0.06) -- (4,-0.06) node[below,font=\scriptsize] {4};
  \draw (5,0.06) -- (5,-0.06) node[below,font=\scriptsize] {5};
  \draw (6,0.06) -- (6,-0.06) node[below,font=\scriptsize] {6};
  \draw (7,0.06) -- (7,-0.06) node[below,font=\scriptsize] {7};
  \draw (8,0.06) -- (8,-0.06) node[below,font=\scriptsize] {8};
  \draw (9,0.06) -- (9,-0.06) node[below,font=\scriptsize] {9};
  \draw (10,0.06) -- (10,-0.06) node[below,font=\scriptsize] {10};
  \draw (11,0.06) -- (11,-0.06) node[below,font=\scriptsize] {11};
  \draw[Blue,thick] (1,5.380) -- (2,4.779) -- (3,3.972) -- (4,3.033) -- (5,1.998) -- (6,0.886);
  \fill[Blue] (1,5.380) circle (1.2pt);
  \fill[Blue] (2,4.779) circle (1.2pt);
  \fill[Blue] (3,3.972) circle (1.2pt);
  \fill[Blue] (4,3.033) circle (1.2pt);
  \fill[Blue] (5,1.998) circle (1.2pt);
  \fill[Blue] (6,0.886) circle (1.2pt);
  \draw[Blue,dashed] (1,5.443) -- (2,4.813) -- (3,3.995) -- (4,3.051) -- (5,2.012) -- (6,0.897);
  \draw[Blue,thick] (7.8,5.78) -- (8.8,5.78);
  \draw[Blue,dashed] (9.1,5.78) -- (10.1,5.78);
  \node[Blue,anchor=west,font=\scriptsize] at (10.3,5.78) {$t=0.5$};
  \draw[black,thick] (1,5.476) -- (2,5.146) -- (3,4.597) -- (4,3.913) -- (5,3.132) -- (6,2.274) -- (7,1.352) -- (8,0.376);
  \fill[black] (1,5.476) circle (1.2pt);
  \fill[black] (2,5.146) circle (1.2pt);
  \fill[black] (3,4.597) circle (1.2pt);
  \fill[black] (4,3.913) circle (1.2pt);
  \fill[black] (5,3.132) circle (1.2pt);
  \fill[black] (6,2.274) circle (1.2pt);
  \fill[black] (7,1.352) circle (1.2pt);
  \fill[black] (8,0.376) circle (1.2pt);
  \draw[black,dashed] (1,5.700) -- (2,5.274) -- (3,4.686) -- (4,3.981) -- (5,3.187) -- (6,2.319) -- (7,1.391) -- (8,0.410);
  \draw[black,thick] (7.8,5.36) -- (8.8,5.36);
  \draw[black,dashed] (9.1,5.36) -- (10.1,5.36);
  \node[black,anchor=west,font=\scriptsize] at (10.3,5.36) {$t=1$};
  \draw[Blue!60!black,thick] (1,5.500) -- (2,5.363) -- (3,5.134) -- (4,4.731) -- (5,4.218) -- (6,3.622) -- (7,2.959) -- (8,2.241) -- (9,1.475) -- (10,0.666);
  \fill[Blue!60!black] (1,5.500) circle (1.2pt);
  \fill[Blue!60!black] (2,5.363) circle (1.2pt);
  \fill[Blue!60!black] (3,5.134) circle (1.2pt);
  \fill[Blue!60!black] (4,4.731) circle (1.2pt);
  \fill[Blue!60!black] (5,4.218) circle (1.2pt);
  \fill[Blue!60!black] (6,3.622) circle (1.2pt);
  \fill[Blue!60!black] (7,2.959) circle (1.2pt);
  \fill[Blue!60!black] (8,2.241) circle (1.2pt);
  \fill[Blue!60!black] (9,1.475) circle (1.2pt);
  \fill[Blue!60!black] (10,0.666) circle (1.2pt);
  \draw[Blue!60!black,dashed] (1,6.114) -- (2,5.862) -- (3,5.478) -- (4,4.995) -- (5,4.432) -- (6,3.802) -- (7,3.114) -- (8,2.377) -- (9,1.596) -- (10,0.775);
  \draw[Blue!60!black,thick] (7.8,4.94) -- (8.8,4.94);
  \draw[Blue!60!black,dashed] (9.1,4.94) -- (10.1,4.94);
  \node[Blue!60!black,anchor=west,font=\scriptsize] at (10.3,4.94) {$t=2$};
\end{tikzpicture}

\caption{\vetnew{The harmonic chain of twelve sites, $\omega=\lambda=1$: the exact
$\abs{\sigma(T_tf,g)}$ (solid, maximized over $f$ at the site $0$ and $g$ at the site $y$ with
coordinates of modulus at most $1$) and the entrywise bound \eqref{eq:harmonicbound} (dashed),
against $y$ on a logarithmic scale, for three times. The bound shares the decay of the exact
quantity and lies above it by a factor that grows with $t$.}}
\label{fig:harmonic}
\end{figure}

\begin{proposition}[witness]\label{prop:harmonicwitness}
\vetnew{For $\lambda>0$ and $n\ge2$, with $f=\delta_{(0,p)}$ and $g=\delta_{(1,p)}$, there is
$t>0$ with $\comm{\alpha_t(W(f))}{W(g)}\ne0$: the momenta at neighbouring sites do not commute
after evolution.}
\end{proposition}
\begin{proof}
\vetnew{$\phi(t)=\sigma(T_tf,g)$ has $\phi(0)=0$ and
$\phi'(0)=\sigma(Kf,g)=(Kf)^q_1=-A_{10}=\lambda>0$, so $0<\phi(t)<2\pi$ for small $t>0$, and the
norm $2\abs{\sin(\phi(t)/2)}$ of Theorem~\ref{thm:weylcommutator} is nonzero. For the position
pair $\delta_{(0,q)},\delta_{(1,q)}$ the first two orders vanish ($Kf$ and $K^2f$ pair with $g$
through momentum coordinates only, and $\sigma(Kf,g)=\sigma(K^2f,g)=0$) and
$\phi(t)=-\lambda t^3/6+O(t^4)$: the commutator appears at order $t^3$.}
\end{proof}

\subsection*{Notes and references}

\vetnew{\emph{The source.} The harmonic bound is Theorem~3.1 of \cite{nrss2009}: for the chain
\eqref{eq:chainH} on a periodic cube (their Hamiltonian without the factor $\tfrac12$, so their
time is half of ours; their chain is periodic where ours has open ends, and in $\nu$ dimensions
their constant carries $4\sum_j\lambda_j$), with $\supp f\subseteq X$, $\supp g\subseteq Y$ and any $\mu>0$,
\begin{gather*}
\norm{\comm{\tau_t(W(f))}{W(g)}}\le C\norm f_\infty\norm g_\infty\sum_{x\in X,\,y\in Y}
\exp\bigl(-\mu\,(d(x,y)-c_{\omega,\lambda}\max(2/\mu,e^{\mu/2+1})\abs t)\bigr),\\
c_{\omega,\lambda}=(\omega^2+4\lambda)^{1/2},
\end{gather*}
their eq.~(3.10), with $C=2+c_{\omega,\lambda}e^{\mu/2}+c_{\omega,\lambda}^{-1}$.

Their proof has the
same first two steps as ours, their Lemma~3.4 ($\tau_t(W(f))=W(f_t)$, obtained there by a
normal-mode computation) and the commutator identity \eqref{eq:weylcomm} (their p.~12 and
eq.~(3.21)); the third step differs: they bound the propagator kernel by a Fourier sum and a
power-series tail with an estimate of the factorials (their Lemmas~3.5 and 3.7), while
Theorem~\ref{thm:entrywise} bounds it entrywise by $e^{\abs t\abs K}$ and Chapter~\ref{sec:closed}'s
tail. Theorem~\ref{thm:harmonicbound} with Corollary~\ref{cor:harmonicclosed} and the paragraph
after it gives their
statement in shape: exponential decay in the distance with a linear light cone, the constants
depending on $\omega$ and $\lambda$ only.

The velocities are not the same number: in their time
unit and at the same $\mu$, our cone (the paragraph after Corollary~\ref{cor:harmonicclosed}) moves at
$2e^\mu c_{\omega,\lambda}/\mu$ after the balanced rescaling, theirs at
$c_{\omega,\lambda}\max(2/\mu,e^{\mu/2+1})$; the two agree at $\mu=2$ and ours is the larger
elsewhere, so their cone is the narrower (at the best $\mu$, $3.6\,c_{\omega,\lambda}$ against
our $5.4\,c_{\omega,\lambda}$); our prefactor has the simpler form.

\emph{The mechanism.} That the dynamics of a quadratic Hamiltonian is a symplectic flow on the
Weyl labels is classical. The bounded-density bounds
for interacting bosons are a different, state-dependent subject and are not treated in these
notes.}

\paragraph{Checks.} {\small \vetnew{The generator condition, symplecticity, the group law, the
Weyl commutator norm, the entrywise bound and the power bound, the chain of inequalities \eqref{eq:harmonicbound} on a chain of
eight sites, the light cone, the row sums before and after the rescaling with the balancing
exponent, the momentum-pair witness and the position pair's $t^3$ onset, the sharper cone tail
and the time normalisation against the source are computed by \texttt{checks/m5\_check.py}; the
numbers of Example~\ref{ex:harmonicchain} (twelve sites) and Figure~\ref{fig:harmonic} by
\texttt{checks/m5\_figure.py}.}}
